\chapter{KẾT LUẬN VÀ HƯỚNG PHÁT TRIỂN}

\section{Kết luận}

Khóa luận đã nghiên cứu, thiết kế và hiện thực thành công một kiến trúc System-on-Chip (SoC) dạng mô-đun dựa trên tập lệnh mở RISC-V, tích hợp các bộ điều khiển ngoại vi truyền thông tiêu chuẩn (UART, SPI, I\textsuperscript{2}C) kết hợp với cơ chế truyền dữ liệu Double Data Rate (DDR). Bằng cách khai thác hiệu quả cả hai cạnh (cạnh lên và cạnh xuống) của tín hiệu xung clock, thiết kế đề xuất đã giúp tăng gấp đôi băng thông truyền dữ liệu trên kênh giao tiếp mà không cần gia tăng tần số clock hệ thống, từ đó hạn chế rủi ro về tiêu thụ năng lượng và nhiễu điện từ (EMI). Đồng thời, cấu trúc mô-đun hóa linh hoạt tạo điều kiện thuận lợi cho việc kiểm chứng, tái sử dụng Intellectual Property (IP) và mở rộng quy mô hệ thống.

Hệ thống được mô tả bằng Verilog HDL và tổng hợp trên AMD Vivado 2025.2 cho FPGA Artix-7 \path{xc7a200tfbg676-2L}. Kết quả cho thấy thiết kế sử dụng 1,814 Slice LUTs, 1,378 Flip-Flops, 1 Block RAM Tile, 48 IOB và 1 BUFGCTRL. Trong mô phỏng chức năng ở $100\text{ MHz}$, khối DDR đạt tốc độ $200\text{ Mbps/pin}$. Báo cáo phân tích định thời với ràng buộc $50\text{ MHz}$ ghi nhận $WNS=+11.659\text{ ns}$ và $WHS=+0.043\text{ ns}$, không có endpoint vi phạm; công suất toàn chip được ước tính là $0.182\text{ W}$. Các kết quả này xác nhận tính khả thi của kiến trúc RISC-V SoC tích hợp giao tiếp DDR cho các ứng dụng nhúng và xử lý tại biên.

\section{Những kết quả đạt được}

Thông qua quá trình nghiên cứu và thực hiện, đề tài đã đạt được các kết quả nổi bật sau:
\begin{itemize}
    \item Về kiến trúc: Lập trình và đóng gói thành công kiến trúc RISC-V SoC dạng mô-đun hóa, hỗ trợ chuẩn kết nối bus nội bộ linh hoạt để dễ dàng tích hợp các khối ngoại vi UART, SPI, I\textsuperscript{2}C.
    \item Về kỹ thuật DDR: Tích hợp thành công khối điều khiển truyền dữ liệu theo định hướng DDR, khai thác tối đa hai cạnh xung clock giúp tối ưu hóa hiệu năng truyền dữ liệu trên các đường giao tiếp ngoại vi.
    \item Về hiện thực RTL: Xây dựng bộ mã nguồn RTL bằng Verilog HDL hoàn chỉnh, tuân thủ các quy tắc tổng hợp logic (Synthesizable RTL), đảm bảo tính đóng gói và khả năng tái sử dụng cao.
    \item Về kiểm chứng: Hoàn thành các testbench UART, SPI, I\textsuperscript{2}C, DDR và tích hợp SoC; đồng thời thực hiện tổng hợp, bố trí định tuyến, phân tích timing, tài nguyên và công suất cho FPGA Artix-7 mục tiêu.
    \item Về tính ứng dụng: Cung cấp một nền tảng SoC phần cứng mở, có khả năng tùy biến cấu hình cao, sẵn sàng phục vụ cho các nghiên cứu phát triển hệ thống nhúng tiếp theo.
\end{itemize}

\section{Những hạn chế}

Mặc dù đã đạt được những kết quả khả quan, thiết kế của khóa luận vẫn tồn tại một số điểm hạn chế cần tiếp tục hoàn thiện:
\begin{itemize}
    \item Tối ưu PPA (Power, Performance, Area): Thiết kế hiện tại chưa được tối ưu sâu về định thời (Timing closure) ở các mức tần số cao; mức độ tiêu thụ tài nguyên (LUTs, Flip-Flops) và công suất tiêu thụ động (Dynamic Power) vẫn còn dư địa để tối ưu.
    \item Cơ chế quản lý dữ liệu: Hệ thống chưa tích hợp bộ điều khiển truy xuất bộ nhớ trực tiếp (DMA), dẫn đến việc bộ xử lý trung tâm (CPU) vẫn phải can thiệp vào quá trình truyền nhận dữ liệu của ngoại vi, làm giảm hiệu năng chung của SoC.
    \item Xử lý miền xung clock (Clock Domain Crossing - CDC): Việc đồng bộ hóa dữ liệu giữa miền xung clock hệ thống và miền xung clock DDR ngoại vi còn đơn giản, cần các giải pháp CDC nâng cao để triệt tiêu hoàn toàn rủi ro trạng thái bất định (Metastability).
    \item Triển khai trên mạch thực: Mã nguồn RTL hiện mới được mô phỏng, tổng hợp và phân tích sau hiện thực trên thiết bị mục tiêu \path{xc7a200tfbg676-2L}. Do chưa có bo mạch FPGA vật lý tương thích, thiết kế chưa thể được đóng gói hoàn chỉnh, gán chân theo phần cứng thực, tạo tệp bitstream và nạp trực tiếp lên bo mạch để kiểm chứng hoạt động.
    \item Môi trường thực nghiệm: Đánh giá hiện dựa trên mô phỏng RTL và báo cáo Implementation của Vivado, chưa đo dạng sóng vật lý bằng Logic Analyzer hoặc Oscilloscope trong điều kiện môi trường công nghiệp.
\end{itemize}

\section{Hướng phát triển trong tương lai}

Để hoàn thiện và nâng cao tính thương mại hóa của thiết kế, các hướng nghiên cứu tiếp theo sẽ tập trung vào:
\begin{itemize}
    \item Tích hợp khối DMA và nâng cấp Bus nội bộ: Tích hợp bộ điều khiển DMA đa kênh và chuyển đổi chuẩn kết nối bus sang các giao thức công nghiệp chuẩn như AMBA AXI4/AHB-Lite nhằm giảm tải tối đa cho CPU.
    \item Tối ưu hóa thiết kế vi mạch: Áp dụng các kỹ thuật thiết kế công suất thấp (Low-power design techniques) như Power Gating, Clock Gating, đồng thời tối ưu đường truyền tới các ô nguyên thủy (Primitives) DDR chuyên dụng trên phần cứng.
    \item Thực hiện quy trình thiết kế ASIC (ASIC Flow): Mở rộng từ FPGA sang quy trình vi mạch ASIC hoàn chỉnh, bao gồm tổng hợp logic với thư viện chuẩn (Standard Cell Library), phân tích định thời tĩnh (STA), thiết kế vật lý (Physical Design - P\&R) và đánh giá công suất tiêu thụ sau đóng gói layout.
    \item Tích hợp bộ gia tốc phần cứng: Mở rộng SoC bằng cách thêm các bộ gia tốc mã hóa/giải mã dữ liệu hoặc bộ gia tốc tính toán AI nhẹ tại biên (TinyML) kết nối qua giao tiếp DDR.
\end{itemize}
